\documentclass[10pt,a4paper]{amsart}
\usepackage[margin=19mm]{geometry}
\usepackage{amsmath,amssymb,mathtools}
\usepackage[scr=rsfs,cal=euler]{mathalfa}
\usepackage{xfrac}
\usepackage[unicode,hidelinks,bookmarks=false]{hyperref}
\newcommand{\ie}{i.e.\ }
\newcommand{\cf}{cf.\ }
\newcommand{\resp}{respectively\ }
\newcommand{\Subsection}{\subsection}
\providecommand{\nobreakdash}{\nobreak\hbox{-}}
\setlength{\emergencystretch}{2em}
\multlinegap0pt
\usepackage{bm}
\usepackage{bbm}
\usepackage{dsfont}
\usepackage{xspace}
\usepackage{cancel}
\usepackage{mathtools}
\usepackage{amsthm}
\makeatletter
\@ifundefined{theorem}{\newtheorem{theorem}{Theorem}}{}
\@ifundefined{proposition}{\newtheorem{proposition}[theorem]{Proposition}}{}
\makeatother
\newcommand{\circM}{\operatorname{circ}}
\title[An operator algebraic approach for generalized Cardano polyn]{An operator algebraic approach for generalized Cardano polynomials}
\author{Leonard Mada, Maria Anastasia Jivulescu}
\date{Source: arXiv:2602.03532v1 (CC BY 4.0)}
\begin{document}
\maketitle

\noindent\emph{Source and licence.} An operator algebraic approach for generalized Cardano polynomials. Version arXiv:2602.03532v1 (CC BY 4.0), distributed under the Creative Commons Attribution 4.0 International licence, as recorded in the arXiv licence metadata for this exact version and shown on the abstract page https://arxiv.org/abs/2602.03532v1. Full rights notice: licenses/arxiv-2602.03532-CC-BY-4.0.txt.\par\smallskip

\noindent\textit{Editorial scope.} Counted unit: the Proposition whose source label is \texttt{th-X} in the article (source0.tex lines 503--510, source label \texttt{th-X}), delivered together with the article's own complete original proof of it (source0.tex lines 511--526, one \texttt{proof} environment of 16 lines). No mathematics is written, repaired or re-derived: statement and proof are the article's own text line for line. Required context retained uncounted: defp-q (lines 256--257); gen-eq (lines 300--313). Every definition, display and label of those lines is retained unchanged. Omitted from this package and not redistributed: 1--255, 258--299, 314--502, 527--733. Nothing inside a retained excerpt is omitted or altered: every retained body is delivered exactly as the article's lines are written, so no omission receipt of any kind accompanies this package. Citations: the retained excerpts carry no citation command at all, so no bibliography entry of the article is transcribed and no reference is redistributed. Reference adaptation: none was needed and none was made; every reference inside the delivered unit resolves inside it, so no unresolved reference or citation is produced. Printed slips kept literal: no printed word, symbol, hypothesis or number of the counted statement or of its proof was altered. Layout and class: the article's own class is \texttt{amsart} with packages including amsmath, amsthm, latexsym, amsfonts, amssymb, xcolor, graphicx, float, graphbox, textcomp, hyperref, geometry, array, comment; the wrapper is the lane's pinned \texttt{amsart} editorial wrapper at 10pt on A4, which restates the theorem-like environments the retained text needs and adds the article's own macro definitions that the retained text needs (\texttt{\textbackslash circM}), with extra wrapper packages (bm, bbm, dsfont, xspace, cancel, mathtools). The delivered numbering is the unit's own. Assets: no retained span contains a figure environment, a graphics-inclusion command, a PDF-page inclusion, a table or a TikZ picture; no image file is copied into this package, the package declares no asset dependency and no image file is redistributed with it. Licence: version arXiv:2602.03532v1 of this article is distributed under the Creative Commons Attribution 4.0 International licence, as recorded in the arXiv licence metadata for this exact version and shown on the abstract page https://arxiv.org/abs/2602.03532v1. A full rights notice is delivered at licenses/arxiv-2602.03532-CC-BY-4.0.txt. Characters: every retained excerpt is pure ASCII, so the delivered engine, XeLaTeX with its default Latin Modern text font, reports no missing character.
\begin{equation}\label{defp-q}
    p:=\sqrt[n]{d+\sqrt{d^2-c^n}} ,\,\, q:=\sqrt[n]{d-\sqrt{d^2-c^n}},\end{equation}

\begin{theorem} Consider $n,$ an odd natural number ($n=2m+1 $), and the real parameters $p$ and $q,$ as given by \eqref{defp-q}.
 The n-th order polynomial equation whose solutions are $x[j]=p\omega^j+q\omega^{-j}, j\in\{0,\pm1,\ldots, \pm [(n-1)/2]\}$ is
 \begin{eqnarray}
x^{2m+1}=2d+\sum\limits_{j=0}^{m-1}B_{m,j}c^{m-j}x^{2j+1}, \text{ where }\label{gen-eq}
\\ B_{m,j}=(-1)^{m-1-j}\frac{2m+1}{2j+1}\begin{pmatrix}
        m+j\\2j
    \end{pmatrix}, 0\leq j\leq m-1.
    \label{B}
\end{eqnarray}

      The polynomial \begin{equation}\label{Car-pol}
          C_{n,c,d}(x)=x^{2m+1}-\sum\limits_{j=0}^{m-1}B_{m,j}c^{m-j}x^{2j+1}-2d,
      \end{equation} associated to the poynomial equation \eqref{gen-eq}  will be called \textit{the n-th order generalized Cardano   polynomial}.
\end{theorem}

\begin{proposition}\label{th-X}
  The Cardano operator $X$ has the followings properties:  
  \begin{itemize}
      \item[a)] It can be written as $X=pX_n+qX_n^{-1}$ ( circular form),  where $X_n$ is the shift operator.
      \item[b)] The operator $X$ has the same spectrum as $W$.
      \item[c)] The operator $X$ satisfies the operator equation $C_{n,c,d}(X)=0$, where $C_{n,c,d}(X)$ is the n-th order Cardano- polynomial for  given real parameters $c$ and $d$, as given by eq.\eqref{gen-eq}.
      \end{itemize}
\end{proposition}

\begin{proof}
    a) Given that $F_n^+Z_nF_n=X_n$ and $F_n^+Z_n^{-1}F_n=X_n^{-1},$ it follows the the operator $X $ can be written as $X=pX_n+qX_n^{-1}$. Given the structure of the operators $X_n$ and $X_n^{-1},$ for which every row is a cycle shift of the first, we conclude that $X$ is a circulant matrix of the 
   form  $X=\circM(0,q,0,\ldots, 0,p)$.
    %\begin{center}
     %   $X=\begin{pmatrix}
      %  0&q&0&\ldots &0&p\\
       % p&0&q&\ldots&0&0\\
        %\vdots&\vdots&\vdots &\ldots &\vdots &\vdots\\
        %0&0&\ldots &p&0&q\\
        %q&0&\ldots &0&p&0
    %\end{pmatrix}$
    %\end{center}
\\b) Given its circular structure as well as the fact that is unitarily equivalent to $W$, the operator $X$ is diagonalizable in the Fourier eigenbasis and its eigenvalues are $\lambda_{\tilde{j}}=p\omega^{j}+q\omega^{-j}.$
\\c) As $F_n$ is a unitary matrix, it holds that $X^k=F_n^+W^kF_n=F_n^+ \text {diag }[\lambda^k[j]]F_n$. \\ We have that 
$C_n(X)=F_n^+C_n(W)F_n=0$.
\end{proof}

\end{document}
